% from main.tex
%%%%%%%%%%  scan idx 384  |  folio 371  %%%%%%%%%%
\textbf{Lemme 1.5.2.} Soient \(A\) un anneau, \(B\) une \(A\)-algèbre de présentation finie, \(M\) un \(B\)-module de présentation finie. On suppose \(B\) et \(M\) plats sur \(A\), et que pour tout corps résiduel \(k\) de \(A\), \(M\otimes_A k\) est de dimension projective \(\leq n\) sur \(B\otimes_A k\) (ce qui est le cas si \(B=A[t_0,\ldots,t_r]\)). Sous ces conditions, \(M\) admet une résolution de longueur \(n\) par des \(B\)-modules projectifs de type fini :
\[
  0\longrightarrow L_n\longrightarrow\cdots\longrightarrow L_0\longrightarrow M\longrightarrow 0.
\]
Ce lemme résulte immédiatement de EGA IV 12.3.2.

\textbf{Proposition 1.6.} Soient \(S\) un schéma quasi-compact, \(E\) un \(\mathcal O_S\)-Module localement libre de rang \(r+1\), \(B\) son algèbre symétrique, avec sa graduation habituelle, \(M\) un \(B\)-Module gradué de présentation finie, plat sur \(S\). Alors il existe une résolution graduée :
\[
  0\longrightarrow L_{r+1}\longrightarrow\cdots\longrightarrow L_0\longrightarrow M\longrightarrow 0,
\]
les \(L_i\) étant des \(B\)-Modules gradués qui sont localement libres de type fini ( en tant que \(B\)-Modules gradués).

Montrons d'abord que pour tout \(B\)-Module gradué \(M\), satisfaisant aux conditions de la proposition, il existe un homomorphisme gradué surjectif~: \(L\to M\), où \(L\) est un \(B\)-Module localement libre de type fini.

Comme \(M\) est de type fini, il existe un entier \(N\) tel que \(M\) soit engendré sur \(B\) par ses composantes de degré \(\leq N\). On pose alors :
\[
  L=\bigoplus_{i\leq N} M_i\otimes_{\mathcal O_S}B(-i);
\]
l'homomorphisme de \(L\) dans \(M\) se définit grâce à (1.2.1), et il est clair qu'il surjectif.

D'après 1.5.2, \(M\) admet une résolution (non graduée) par des \(B\)-Modules localement libres de type fini sur chaque ouvert affine. Par suite, si \(M\) a une
